\section{Introduction}
\label{sec:introduction}
Brewery Plat-Eau is an optimisation project about choosing how to produce beer at minimum cost. The model decides how much barley, malt, malt extract, hops, and yeast to use for a requested beer volume and colour. Later tasks also require minimum ingredient quality.

The continuous decision variables are ingredient quantities. Two binary variables represent whether the malting and mashing processes are activated, because these processes have fixed costs. The report is organised as a direct answer to the assignment tasks: first the common brewing assumptions and optimisation model are defined, and then Tasks 1--7 are answered in order.
